
\documentclass{article}
\usepackage{colortbl}
\usepackage{makecell}
\usepackage{multirow}
\usepackage{supertabular}

\begin{document}

\newcounter{utterance}

\centering \large Interaction Transcript for game `wordle\_withclue', experiment `high\_frequency\_words\_clue\_no\_critic', episode 2 with reflexion{-}Qwen3.5{-}27B.
\vspace{24pt}

{ \footnotesize  \setcounter{utterance}{1}
\setlength{\tabcolsep}{0pt}
\begin{supertabular}{c@{$\;$}|p{.15\linewidth}@{}p{.15\linewidth}p{.15\linewidth}p{.15\linewidth}p{.15\linewidth}p{.15\linewidth}}
    \# & \multicolumn{2}{c}{Player} && \multicolumn{2}{c}{Game Master} \\
    \hline

    \theutterance \stepcounter{utterance}  
    & & & \multicolumn{4}{p{0.6\linewidth}}{
        \cellcolor[rgb]{0.9,0.9,0.9}{
            \makecell[{{p{\linewidth}}}]{
                \texttt{\tiny{[P1$\langle$GM]}}
                \texttt{You have attempted this same task before and did not succeed. These are the plans you wrote after those attempts.} \\
\texttt{Trial \#1: New plan: The previous attempt was aborted because the response violated the strict formatting rules by including a long, rambling explanation that failed to conclude with a valid "guess:" tag and a five{-}letter word. The explanation drifted into a list of rejected words without ever committing to a final guess. Next time, I will strictly adhere to the required format: a concise "explanation:" followed immediately by "guess: {[}word{]}". I will ensure the guess is a valid five{-}letter English word that respects the feedback (excluding red letters 'd', 'a', 'n', 'c' and placing 'e' in a position other than the 5th) and fits the clue. For example, I will select a word like "reels" (if 'r', 'l', 's' are valid) or "beeps" (if 'b', 'p', 's' are valid), ensuring the output contains no extra text or incomplete thoughts.} \\
\\ 
\texttt{Trial \#2: New plan: The previous attempt was aborted because the response failed to provide a complete "guess:" tag and a valid five{-}letter word, instead trailing off with an incomplete list of possibilities. I violated the rule "Do not generate any other text" by including a rambling, unfinished thought process. Next time, I will strictly adhere to the format: a concise "explanation:" followed immediately by "guess: {[}word{]}". I will ensure the guess is a valid five{-}letter English word that respects the feedback (excluding red letters 'd', 'a', 'n', 'c' and placing 'e' in the 5th position as indicated by the green feedback) and fits the clue. For example, I will select a word like "voice" or "scale" and ensure the output ends definitively with the guess, containing no extra text or incomplete sentences.} \\
\\ 
\texttt{Trial \#3: New plan: The previous attempt was aborted because the response failed to provide a complete "guess:" tag and a valid five{-}letter word, instead trailing off with an incomplete list of possibilities and internal monologue. I violated the rule "Do not generate any other text" by including a rambling, unfinished thought process. Next time, I will strictly adhere to the format: a concise "explanation:" followed immediately by "guess: {[}word{]}". I will ensure the guess is a valid five{-}letter English word that respects the feedback (excluding red letters 'l', 's', 'g', 'n' and placing 'e' in the 4th position as indicated by the green feedback, while placing 'r' in a position other than 1st or 2nd) and fits the clue. For example, I will select a word like "three" or "shred" (if 's' wasn't red, but it was, so "three" is better) and ensure the output ends definitively with the guess, containing no extra text or incomplete sentences.} \\
\\ 
\\ 
\texttt{You are a language wizard who likes to guess words by using the given rules.} \\
\\ 
\texttt{Welcome to Wordle! You have six attempts to guess the target word, a valid English word of five lowercase letters (a{-}z). Please use the tags "explanation:" and "guess:" to provide a concise explanation for each guess.} \\
\\ 
\texttt{To help you make an informed guess, you will receive a clue for the word, such as} \\
\texttt{clue: snowy white.} \\
\\ 
\texttt{Here is an example guess based on the clue:} \\
\texttt{explanation: In the fairy tail Snow White, the girl is killed because she eats a poisoned apple. And the word apple has 5 letters.} \\
\texttt{guess: apple} \\
\\ 
\texttt{After each guess, your answer will be validated, and you will receive feedback indicating which letters are correct (green), which letters are correct but in the wrong position (yellow), and which letters are incorrect (red). This feedback can be useful in determining which letters to include or exclude in your next guess.} \\
\\ 
\texttt{For example, the feedback for "apple" might be:} \\
\texttt{guess\_feedback: a<yellow> p<yellow> p<green> l<yellow> e<red>} \\
\\ 
\texttt{The explanation should contain details about how both the guess\_feedback and clue work together in tandem to formulate a new guess.} \\
\\ 
\texttt{Do not generate any other text.} \\
\\ 
\texttt{Let's begin with your first guess.} \\
\\ 
\texttt{clue: "...called for his fiddlers \_\_\_"} \\
            }
        }
    }
     \\ \\

    \theutterance \stepcounter{utterance}  
    & \multicolumn{4}{p{0.6\linewidth}}{
        \cellcolor[rgb]{0.561,0.737,0.561}{
            \makecell[{{p{\linewidth}}}]{
                \texttt{\tiny{[P1$\rangle$GM]}}
                \texttt{explanation: The clue "...called for his fiddlers \_\_\_" is a famous line from the nursery rhyme "The Grand Old Duke of York," which continues with "to play." The word "play" is four letters, but the rhyme often implies a specific action or instrument. However, a more direct fit for a 5{-}letter word in this context, often associated with fiddlers in riddles or specific song lyrics, might be "music" or "dance". Let's reconsider the specific phrase. The line is "He called for his fiddlers to play." Wait, "to play" is two words. Is there a 5{-}letter word? Perhaps the clue refers to the word "fiddle" itself? No, that's 6. Let's look at the phrase "called for his fiddlers to \_\_\_". The missing word is likely "play". But "play" is 4 letters. Maybe the clue is "The fiddlers \_\_\_"? Or perhaps the word is "tune"? No. Let's try a common 5{-}letter word associated with fiddlers: "music". Or maybe the clue is a misquote of "The fiddlers on the roof"? No. Let's try "dance". Fiddlers often play for dancing. "Dance" is 5 letters. Another possibility: "reel" is 4. "Jig" is 3. "Song" is 4. "Tune" is 4. "Band" is 4. "Orch" is not a word. "Folk" is 4. "Play" is 4. Is it possible the word is "music"? "He called for his fiddlers music"? No. Let's re{-}read the clue carefully: "...called for his fiddlers \_\_\_". Maybe the word is "away"? No. "Here"? No. "Now"? No. "Soon"? No. "Fast"? No. "Hard"? No. "Loud"? No. "Soft"? No. "Slow"? No. "Stop"? No. "Start"? 5 letters. "Start" fits. "He called for his fiddlers start"? No. "Play" is the most logical verb, but it's 4 letters. Is there a 5{-}letter synonym? "Perform"? 7. "Entertain"? 9. "Amuse"? 5. "He called for his fiddlers amuse"? No. Maybe the word is "fiddle"? No, 6. "Fiddlers" is the subject. The object is missing. "He called for his fiddlers {[}to{]} play". Maybe the word is "tunes"? 5 letters. "He called for his fiddlers tunes"? No. "He called for his fiddlers to play tunes". Maybe the word is "music". "He called for his fidd} \\
            }
        }
    }
    & & \\ \\

    \theutterance \stepcounter{utterance}  
    & & & \multicolumn{2}{p{0.3\linewidth}}{
        \cellcolor[rgb]{0.9,0.9,0.9}{
            \makecell[{{p{\linewidth}}}]{
                \texttt{\tiny{[GM$|$GM]}}
                \texttt{Error: The response should contain the 'guess:' keyword exactly once.} \\
            }
        }
    }
    & & \\ \\

    \theutterance \stepcounter{utterance}  
    & & & \multicolumn{2}{p{0.3\linewidth}}{
        \cellcolor[rgb]{0.9,0.9,0.9}{
            \makecell[{{p{\linewidth}}}]{
                \texttt{\tiny{[GM$|$GM]}}
                \texttt{game\_result = ABORT} \\
            }
        }
    }
    & & \\ \\

    \theutterance \stepcounter{utterance}  
    & & \multicolumn{4}{p{0.6\linewidth}}{
        \cellcolor[rgb]{0.95,0.95,0.95}{
            \makecell[{{p{\linewidth}}}]{
                \texttt{\tiny{[P1$\leftrightarrow$P]}}
                \texttt{Post{-}episode Reflexion (not sent to GM):} \\
\texttt{New plan: The previous attempt was aborted because the response violated the strict formatting rule "Do not generate any other text" by including a long, rambling internal monologue that never concluded with a valid "guess:" tag. Instead of outputting a thought process, I must directly provide the required tags. Next time, I will immediately generate a concise "explanation:" that connects the clue ("...called for his fiddlers \_\_\_") to a specific five{-}letter word (likely "tunes" or "music" based on the context of fiddlers playing), followed immediately by "guess: {[}word{]}". I will ensure the guess is a valid five{-}letter English word and that the response contains absolutely no extra text, incomplete sentences, or visible deliberation.} \\
            }
        }
    }
    & \\ \\

\end{supertabular}
}

\end{document}
